\section{Related Work}
\label{sec:related}

\para{Self-correction and self-critique in LLMs.}
The idea that LLMs can improve their own outputs through self-reflection has driven systems such as Self-Refine \cite{madaan2023selfrefine}, which iteratively generates feedback and revisions, and Reflexion \cite{shinn2023reflexion}, which uses verbal reinforcement learning for agent tasks. However, \citet{huang2024selfcorrect} demonstrated that intrinsic self-correction---without external oracle feedback---degrades performance on reasoning benchmarks including \gsmk and CommonSenseQA. They showed that models are more likely to change correct answers to incorrect ones than vice versa. Our work extends this line of inquiry from behavioral outcomes to internal mechanisms, asking \emph{why} self-correction fails at the representational level.

\para{Probing internal representations.}
A substantial body of work has used linear probes to study what information is encoded in LLM hidden states. \citet{marks2024geometry} demonstrated that truth values are linearly separable in the residual stream and that mass-mean probe directions are causally implicated in model predictions. \citet{burns2023discovering} proposed Contrast Consistent Search (CCS) for unsupervised extraction of truth directions. \citet{garner2025probing} showed that models internally represent correct arithmetic answers even when outputting incorrect ones, achieving over 90\% error detection accuracy from residual stream probes. \citet{david2025temporal} found that linear probes can predict chain-of-thought correctness from as few as 4 generated tokens. These studies establish that residual stream activations carry rich information about reasoning quality, but none has examined how self-critique affects this information.

\para{Representation geometry and engineering.}
\citet{zou2023representation} introduced Representation Engineering (RepE), using Linear Artificial Tomography to extract and manipulate concept directions in hidden states. \citet{anderson2026geometry} analyzed reasoning trajectories through representation space, identifying universal geometric signatures across domains including oscillatory step-to-step coherence patterns. \citet{turpin2024unfaithful} showed that chain-of-thought explanations can be unfaithful to actual model reasoning, motivating the study of internal states rather than surface-level outputs. Our work applies geometric analysis tools from this literature to the specific perturbation caused by self-critique, measuring how the representational subspace changes across the critique-revision cycle.

\para{Positioning of our work.}
Unlike prior studies that examine either the behavioral effects of self-correction \cite{huang2024selfcorrect,madaan2023selfrefine} or the static structure of representations during initial generation \cite{marks2024geometry,garner2025probing,david2025temporal}, we track how representations \emph{shift} across the generate-critique-revise pipeline. We combine the self-correction evaluation framework of \citet{huang2024selfcorrect} with the probing methodology of \citet{marks2024geometry} and the geometric analysis of \citet{anderson2026geometry} to provide the first mechanistic account of why self-critique fails.
